\scanpage{254}
\mspage{200}

\[
\mathrm{HOT}^{\unsure{\mathrm{lc}|\mathrm{le}}}_W(S)\times\mathrm{HOT}^{\mathrm{lc}}_W(S)\to\mathrm{HOT}^{\mathrm{lc}}_W(S)
\]
lui-même induit par le foncteur \inserted{produit} $(X,Y)\mapsto X\times_S Y$ dans $\underline{\mathrm{Fib}}_W S$.

Si $X$ est quelconque dans $\mathrm{Cat}_{/S}$, le foncteur
\[
\mathrm{mult}_{X/S}:Y\to X\times_S Y:\ \underline{\mathrm{Liss}}_W S\to\mathrm{Cat}_{/S}
\]
est compatible avec les localiseurs $W_S^d$, car ce foncteur peut s'écrire comme le composé du foncteur correspondant à \uline{val. ds} \deleted{$\mathrm{Cat}_{/X}$} $\underline{\mathrm{Liss}}_W X$, lequel est compatible avec les localiseurs $W_S^d,W_X^d$, et le foncteur ``oubli de $X$'' de $\underline{\mathrm{Liss}}_W X$ (\deleted{ou} $=$ de $\mathrm{Cat}_{/X}$) dans $\mathrm{Cat}_{/S}$, lequel est compatible avec $W_X^d$ et $W_S^d$. On trouve ainsi un foncteur
\[
\left\{\begin{aligned}
\mathrm{mult}_{X/S}:&\quad\mathrm{Hot}_W(S)\to\mathrm{Hot}_W(S)\\
\mathrm{mult}_{X/S}(\mathrm{hot}_{S,W}\deleted{\unsure{\mathrm{hot}}}(Y))&\overset{\mathrm{def}}{=}\mathrm{hot}_{S,W}(X\times_S Y),\ \text{si }Y\ \underline{\text{lisse}}\text{ sur }S^{*}
\end{aligned}\right.
\]
% slanted
\margin{Cette formule est également valable si $Y$ est dans $\underline{\mathrm{Fib}}_W S$, cf $(*)$ p. 202.}
\unsure{puisqu'ici|puisque} \deleted{\unsure{ici}} $\mathrm{HOT}_W(S)$ peut s'obtenir comme localisé tant de $\underline{\mathrm{Liss}}_W S$ que de $\mathrm{Cat}_{/S}$, par $W_S^d$. Quand $X$ est lui-$=$ \inserted{$W$-}lisse sur \deleted{\unsure{sa}} $S$, ce foncteur n'est autre que
\[
\left\{\begin{aligned}
\mathrm{mult}_{\unsure{\xi|\zeta}}:&\quad\unsure{\eta|y}\longmapsto\xi\times_S\eta,\\
&\text{où }\xi=\mathrm{hot}_{S,W}(X).
\end{aligned}\right.
\]
Mais il n'en est pas de $=$ si $X$ n'est pas lisse, i.e. il faut distinguer dans ce cas, même \inserted{pour $Y$ \unsure{$W$-|lui-$=$ }lisse sur $S$}
\[
\mathrm{hot}_{S,W}(X\times_S Y)\quad\text{et}\quad\mathrm{hot}_{S,W}(X)\times_S\mathrm{hot}_{S,W}(Y).
\]
\deleted{\unsure{Pour $X$ et $Y$ qui}} \inserted{Le deuxième terme peut} s'expliciter comme $\mathrm{hot}_{S,W}(\unsure{\Phi|\mathscr{D}}(X)\times_S Y))$ qui reçoit $\mathrm{hot}_{S,W}(X\times_S Y)$, i.e.
% archive sheet number: 253
